\documentclass[10.5pt]{article}
\usepackage{fontspec}\defaultfontfeatures{Ligatures=TeX}
\usepackage[letterpaper,margin=0.6in]{geometry}
\usepackage{amsmath,bm,xcolor,longtable,array,tikz,cancel}
\usetikzlibrary{arrows.meta,patterns}
\setlength{\parindent}{0pt}\setlength{\parskip}{4pt}
\definecolor{pblue}{HTML}{1971C2}\definecolor{pred}{HTML}{E03131}\definecolor{pgreen}{HTML}{2F9E44}\definecolor{ppurple}{HTML}{7048E8}\definecolor{pgray}{HTML}{868E96}
\newfontfamily\mono{TeX Gyre Cursor}
\newcommand{\mtag}[2]{\colorbox{#1}{\textcolor{white}{\textsf{\fontsize{6.5}{7}\selectfont\bfseries #2}}}}
\newcommand{\clock}[1]{{\mono\fontsize{7}{8}\selectfont #1}\ }
\newcommand{\step}[4]{\par\noindent\makebox[0.32in][l]{\textbf{\small #1}}\mtag{#2}{#3}\ {\small #4}\par\vspace{2pt}}
\newcommand{\pic}[1]{\par\hspace{0.32in}\begin{tikzpicture}[scale=0.42,every node/.style={font=\scriptsize}]#1\end{tikzpicture}\par\vspace{3pt}}
\newcommand{\CHECK}[3]{\par\noindent\makebox[0.32in][l]{\textbf{\small #1}}\fcolorbox{black}{black!4}{\parbox{2.55in}{\textsf{\fontsize{6.5}{8}\selectfont\bfseries CHECK}\ {\fontsize{7.5}{9}\selectfont #2}\par{\fontsize{7.5}{9}\selectfont\textbf{Answer:} #3}\par{\fontsize{7}{8.5}\selectfont\color{pgray} 4 of 5 right: go on. Fewer: one more example, ask again.}}}\par\vspace{3pt}}
\newcommand{\WATCH}[1]{\par\hspace{0.32in}{\textsf{\fontsize{6.5}{8}\selectfont\bfseries\color{pred} WATCH FOR}\ {\small\itshape\color{pred} #1}}\par}
\newcommand{\bl}[1]{\underline{\textcolor{pgreen}{\textbf{#1}}}}
\newcommand{\pen}[2]{\textcolor{#1}{\rule{8pt}{8pt}}\ #2}
% the floor and the crate, as preset on the board: \scene{push label}
\newcommand{\floor}{\fill[pattern=north east lines,pattern color=pgray] (-3,-0.3) rectangle (3,0);\draw[thick] (-3,0)--(3,0);}
\newcommand{\crate}{\draw[pblue,thick,fill=pblue!8] (-0.8,0) rectangle (0.8,1.4);\node at (0,0.7){12 kg};}
\newcommand{\pushA}[1]{\draw[pblue,very thick,-{Stealth}] (-2.8,0.7)--(-0.85,0.7);\node[pblue,above] at (-1.8,0.7){#1};}
\newcommand{\frc}[3]{\draw[pred,very thick,-{Stealth}] #1--#2;\node[pred] at #3;}
% the friction graph axes (push 0-80 across, friction 0-60 up; 1 unit = 10 N) with the two dashed levels
\newcommand{\axes}{\draw[pblue,thick,-{Stealth}] (0,0)--(8.6,0);\draw[pblue,thick,-{Stealth}] (0,0)--(0,6.6);\node[below] at (4,-1.0){push (N)};\node[rotate=90] at (-1.7,3){friction (N)};
 \foreach \v in {2,4,6,8}{\node[below,pgray] at (\v,0){\v0};}\foreach \v in {2,4,6}{\node[left,pgray] at (0,\v){\v0};}
 \draw[pgray,dashed] (0,4.7)--(8,4.7) node[right]{47};\draw[pgray,dashed] (0,2.94)--(8,2.94) node[right]{29};}
% the two carts and rope; the table rig; the Atwood machine (as preset on the board)
\newcommand{\carts}{\fill[pattern=north east lines,pattern color=pgray] (-0.2,-0.25) rectangle (7.4,0);\draw[thick] (-0.2,0)--(7.4,0);
 \draw[pblue,thick,fill=pblue!8] (0,0) rectangle (1.4,1);\node at (0.7,0.5){2.0 kg};\draw[pblue,thick,fill=pblue!8] (3,0) rectangle (5,1);\node at (4,0.5){4.0 kg};
 \draw[brown,thick] (1.4,0.5)--(3,0.5);\draw[pblue,very thick,-{Stealth}] (5,0.5)--(7,0.5) node[above left]{12 N};}
\newcommand{\tbl}{\draw[thick] (0,2)--(5,2)--(5,-0.6);\draw[fill=black!20] (5.25,2.2) circle (0.22);
 \draw[pblue,thick,fill=pblue!8] (1,2) rectangle (2.8,2.9);\node at (1.9,2.45){5.0 kg};\draw[brown,thick] (2.8,2.42)--(5.25,2.42);\draw[brown,thick] (5.47,2.2)--(5.47,1);
 \draw[pblue,thick,fill=pblue!8] (5.07,0.2) rectangle (5.87,1);\node[right] at (5.9,0.6){2.0 kg};}
\newcommand{\atw}{\draw[thick] (0.4,4)--(3.6,4);\draw (2,4)--(2,3.6);\draw[fill=black!20] (2,3.1) circle (0.55);\draw[brown,thick] (1.45,3.1)--(1.45,1.7);\draw[brown,thick] (2.55,3.1)--(2.55,1.1);
 \draw[pblue,thick,fill=pblue!8] (1.05,1) rectangle (1.85,1.7);\node at (1.45,1.35){4.0};\draw[pblue,thick,fill=pblue!8] (2.15,0.2) rectangle (2.95,1.1);\node at (2.55,0.65){6.0};}
\pagestyle{plain}
\begin{document}
\begin{center}{\Large\bfseries Forces, Day 4 board script: Connected Objects}\\[2pt]{\normalsize Honors \textbullet{} 50 minutes \textbullet{} \texttt{physics/whiteboard/index.html\#lesson=systems}}\end{center}
\noindent\textbf{Pens.} \pen{pblue}{objects, axes, motion} \quad \pen{pred}{forces} \quad \pen{pgreen}{new words (labels)} \quad \pen{ppurple}{components} \quad \pen{black}{numbers and everything else}\par
\noindent\textbf{Steps.} Numbered by frame: 1.1, 1.2\ldots{} are on frame 1. Each DRAW step has a small picture of what the board should look like when the step is done. Tick the green Words box as you write each label.\par
\noindent\textbf{Numbers all day.} Two carts, 2.0 and 4.0 kg, pulled with 12 N: $a = 2.0$ m/s$^2$, $T = 4.0$ N. A 5.0 kg cart and a 2.0 kg hanging mass: $a = 2.8$ m/s$^2$, $T = 14$ N. Atwood, 4.0 and 6.0 kg: $a = 1.96$ m/s$^2$, $T = 47$ N. The cart with $\mu_k = 0.20$: $f_k = 9.8$ N, $a = 1.4$ m/s$^2$, $T = 16.8$ N. Every tension checks from both objects. None of these are test numbers; do not use 3.0 kg with 1.0 kg, or 5.0 kg with 3.0 kg.\par
\noindent\textbf{Not today.} No pulleys with mass, no more than one pulley, no ramps with a counterweight (that is the Day 8 cargo ramp), no tension in a rope that has mass (Day 6 practice).\par\smallskip
\noindent\fcolorbox{black}{black!3}{\parbox{\dimexpr\linewidth-2\fboxsep}{\small
\textbf{Time checkpoints.} Be starting frame 2 by \textbf{14:00}, frame 3 by \textbf{24:00}, frame 4 by \textbf{31:00}, the exit check by \textbf{38:00}. The checks are never cut; trim around them.\par
\textbf{If you are behind,} cut in this order: (1) step 2.4, the check from the hanging mass: say it in one sentence; (2) step 3.4, the tension from the 6.0 kg side; (3) frame 4's worked tension (keep $a$, drop $T$); (4) exit question b; (5) the summary becomes tonight's homework.\par
\textbf{If you are ahead} (the 4-of-5 checks pass quickly), use these before the exit check, one at a time on whiteboards: \emph{Three carts, 1.0, 2.0 and 3.0 kg, pulled from the front with 18 N. Tension between the back two?} (a = 3.0, T = 3.0 N.) \emph{What hanging mass makes the 5.0 kg cart speed up at 4.9 m/s$^2$?} (5.0 kg: half of g when the masses are equal.) \emph{In the Atwood machine, what happens to T if both masses double?} (It doubles to 94 N; a stays 1.96.) \emph{With $\mu_k = 0.40$, does the 5.0 kg cart move at all?} (Static friction would decide; with $\mu_s \ge 0.40$, 19.6 N cannot beat 19.6 N, so no.)}}\par\medskip
\begin{longtable}{@{}>{\setlength{\parskip}{5pt}}p{3.55in}@{\hspace{0.25in}}p{3.05in}@{}}
\multicolumn{2}{@{}l}{\large\bfseries Frame 0: question and warm-up \hfill{\mono\small 0:00--7:00}}\tabularnewline*[2pt]
\textbf{Essential question.} When objects are tied together, how do we find how fast they speed up and how hard the string pulls?\par
\textbf{Warm-up.} 12 N of net force on 3.0 kg: acceleration? Weight of a 2.0 kg mass? A rope pulls a box: its third-law partner? &
\step{0.1}{black}{SAY}{\clock{0:00}Warm-up is up as they walk in. At 3:00 remove the covers: 4.0 m/s$^2$; 19.6 N; the box pulls back on the rope.}
\step{0.2}{black}{SAY}{\clock{5:00}Read the question. ``Today, two objects and one string. The trick is choosing what to look at.''} \tabularnewline[10pt]
\multicolumn{2}{@{}l}{\large\bfseries Frame 1: systems \hfill{\mono\small 7:00--14:00}}\tabularnewline*[2pt]
Treat the tied objects as one \bl{system}. The string pulls the two objects with equal and \bl{opposite} forces, so tension is \bl{internal} and cancels. Only \bl{external} forces change how the system moves.\par
System: $a$ = (net external force) $\div$ (total mass). Then look at \bl{one} object to find the tension.\par
\textbf{Example.} 2.0 and 4.0 kg carts, 12 N: $a = 12 \div 6.0 = $ \bl{2.0 m/s$^2$}. The back cart: $T = 2.0 \times 2.0 = $ \bl{4.0 N}.\par\smallskip
\textbf{1.} A student says the rope between the two carts pulls with 12 N, because that is the pull. What is wrong? &
\step{1.1}{pred}{DRAW}{Tension on both carts: the rope pulls the back cart forward and the front cart back. Same length. Say: ``Third-law pair, through the rope.'' Tick \textbf{opposite}.}
\pic{\carts\draw[pred,very thick,-{Stealth}] (1.4,0.5)--(2.1,0.5);\draw[pred,very thick,-{Stealth}] (3,0.5)--(2.3,0.5);\node[pred] at (2.2,1){T};}
\step{1.2}{pblue}{DRAW}{A dashed box around both carts and the rope. Write SYSTEM, 6.0 kg. Say: ``Inside the box, the two T arrows cancel. INTERNAL. The 12 N comes from outside: EXTERNAL.'' Tick \textbf{system}, \textbf{internal}, \textbf{external}.}
\pic{\carts\draw[pblue,dashed] (-0.15,-0.1) rectangle (5.15,1.15);\node[pgreen] at (2.5,1.5){SYSTEM: 6.0 kg};}
\tabularnewline
&
\step{1.3}{black}{DRAW}{Under the picture: $a = 12 \div 6.0 = 2.0$ m/s$^2$. Tick \textbf{2.0 m/s$^2$}, then circle the 2.0 kg cart: ``Now ONE object. What pulls it forward?'' (Only T.) $T = 2.0 \times 2.0 = 4.0$ N. Tick \textbf{one}, \textbf{4.0 N}.}
\pic{\node[right] at (0,0){$a = 12 \div 6.0 = \textcolor{pgreen}{\mathbf{2.0\ m/s^2}}$};\node[right] at (0,-1){back cart: $T = 2.0 \times 2.0 = \textcolor{pgreen}{\mathbf{4.0\ N}}$};}
\CHECK{1.4}{The rope between the carts pulls with 12 N?}{No: it only speeds up the 2.0 kg cart, $T = 4.0$ N. The 12 N speeds up both.}
\WATCH{Dividing 12 N by one cart's mass. Ask: which objects does the 12 N have to move?} \tabularnewline[10pt]
\multicolumn{2}{@{}l}{\large\bfseries Frame 2: a table and a hanging mass \hfill{\mono\small 14:00--24:00}}\tabularnewline*[2pt]
The only external force along the string is the hanging mass's \bl{weight}. It has to speed up \bl{both} masses.\par
\textbf{Example.} 5.0 kg cart, frictionless table, 2.0 kg hanging. System: $a = 19.6 \div 7.0 = $ \bl{2.8 m/s$^2$}. Cart alone: $T = 5.0 \times 2.8 = $ \bl{14 N}. Check: $19.6 - T = 2.0 \times 2.8$, so $T = 14$ N. $T$ is \bl{less} than its weight.\par\smallskip
\textbf{2.} A student says the hanging mass falls with $a = g = 9.8$ m/s$^2$. What is wrong? &
\step{2.1}{black!60}{ASK}{``Let go. What makes this move?'' (The hanging weight.) ``What has to speed up?'' (Both.) Tick \textbf{weight}, \textbf{both}.}
\step{2.2}{pred}{DRAW}{$W = 19.6$ N down on the hanging mass; $T$ along the string on both. Say: ``The string bends over the pulley, so we follow the string as if it were straight.''}
\pic{\tbl\draw[pred,very thick,-{Stealth}] (5.47,0.2)--(5.47,-0.7) node[right]{19.6 N};\draw[pred,very thick,-{Stealth}] (2.8,2.42)--(3.8,2.42) node[above]{T};\draw[pred,very thick,-{Stealth}] (5.47,1)--(5.47,1.7);}
\tabularnewline
&
\step{2.3}{black}{DRAW}{In ``Work it here'': $a = 19.6 \div 7.0 = 2.8$ m/s$^2$; cart alone: $T = 5.0 \times 2.8 = 14$ N. Tick \textbf{2.8 m/s$^2$}, \textbf{14 N}.}
\pic{\node[right] at (0,0){$a = 19.6 \div 7.0 = \textcolor{pgreen}{\mathbf{2.8\ m/s^2}}$};\node[right] at (0,-1){cart: $T = 5.0 \times 2.8 = \textcolor{pgreen}{\mathbf{14\ N}}$};}
\step{2.4}{black!60}{ASK}{``Check from the hanging mass.'' ($19.6 - T = 2.0 \times 2.8$, $T = 14$ N.) ``Why is 14 less than 19.6?'' (It speeds up downward, so the net force is down.) Tick \textbf{less}.}
\CHECK{2.5}{The hanging mass falls with $a = g$?}{No: its weight must speed up the cart too, so $a$ is less than $g$: 2.8 m/s$^2$.}
\WATCH{$T = 19.6$ N, ``the string holds the weight.'' Ask: is the hanging mass at rest?} \tabularnewline[10pt]
\multicolumn{2}{@{}l}{\large\bfseries Frame 3: the Atwood machine \hfill{\mono\small 24:00--31:00}}\tabularnewline*[2pt]
The \bl{heavier} side goes down. The net external force is the \bl{difference} of the two weights, and it speeds up the \bl{total} mass.\par
\textbf{Example.} 4.0 and 6.0 kg: $a = (58.8 - 39.2) \div 10 = $ \bl{1.96 m/s$^2$}. The 4.0 kg mass goes up: $T - 39.2 = 4.0 \times 1.96$, so $T = $ \bl{47 N}.\par\smallskip
\textbf{3.} Two 3.0 kg masses hang over a pulley. Find $a$ and $T$. &
\step{3.1}{black!60}{ASK}{``Which way?'' (6.0 kg down.) Draw blue $a$ arrows: down on 6.0, up on 4.0. Tick \textbf{heavier}.}
\step{3.2}{pred}{DRAW}{Weights 58.8 N and 39.2 N down; $T$ up on both, the same $T$. Say: ``One string, one tension.''}
\pic{\atw\draw[pred,very thick,-{Stealth}] (1.25,1)--(1.25,0.2) node[left]{39.2};\draw[pred,very thick,-{Stealth}] (2.75,0.2)--(2.75,-0.6) node[right]{58.8};\draw[pred,very thick,-{Stealth}] (1.25,1.7)--(1.25,2.3) node[left]{T};\draw[pred,very thick,-{Stealth}] (2.75,1.1)--(2.75,1.7) node[right]{T};}
\tabularnewline
&
\step{3.3}{black}{DRAW}{$a = (58.8 - 39.2) \div 10 = 1.96$ m/s$^2$. Tick \textbf{difference}, \textbf{total}, \textbf{1.96 m/s$^2$}.}
\step{3.4}{black}{DRAW}{The 4.0 kg mass alone, up positive: $T - 39.2 = 4.0 \times 1.96$, $T = 47$ N. Check from the 6.0 kg side: $58.8 - T = 6.0 \times 1.96$, $T = 47$ N. Tick \textbf{47 N}.}
\pic{\node[right] at (0,0){$a = (58.8 - 39.2) \div 10 = \textcolor{pgreen}{\mathbf{1.96}}$};\node[right] at (0,-1){$T - 39.2 = 4.0 \times 1.96$, $T = \textcolor{pgreen}{\mathbf{47\ N}}$};}
\CHECK{3.5}{Two 3.0 kg masses over a pulley: $a$ and $T$?}{Equal weights: $a = 0$, $T = 29.4$ N.}
\WATCH{$T = 58.8 - 39.2 = 19.6$ N. Ask: then could the 4.0 kg mass rise?} \tabularnewline[10pt]
\multicolumn{2}{@{}l}{\large\bfseries Frame 4: friction on the table \hfill{\mono\small 31:00--38:00}}\tabularnewline*[2pt]
Friction is one more external force, pointing \bl{against} the motion. It acts only on the table mass, so use \bl{that mass's} $N$.\par
\textbf{Example.} The 5.0 kg cart with $\mu_k = 0.20$: $f_k = 0.20 \times 49 = $ \bl{9.8 N}. $a = (19.6 - 9.8) \div 7.0 = $ \bl{1.4 m/s$^2$}. Cart: $T - 9.8 = 5.0 \times 1.4$, so $T = $ \bl{16.8 N}.\par\smallskip
\textbf{4.} A student finds the friction with $N = (5.0 + 2.0) \times 9.8$. What is wrong? &
\step{4.1}{pred}{DRAW}{Friction on the cart, pointing left, at its bottom face. Tick \textbf{against}. Ask: ``Which mass presses on the table?'' (The cart only.) Tick \textbf{that mass's N}.}
\pic{\tbl\draw[pred,very thick,-{Stealth}] (2.2,2.03)--(1.1,2.03) node[above left]{$f_k$};}
\step{4.2}{black}{DRAW}{$f_k = 0.20 \times 49 = 9.8$ N; $a = (19.6 - 9.8) \div 7.0 = 1.4$ m/s$^2$; cart: $T - 9.8 = 5.0 \times 1.4$, $T = 16.8$ N. Tick \textbf{9.8 N}, \textbf{1.4 m/s$^2$}, \textbf{16.8 N}.}
\CHECK{4.3}{Friction with $N = (5.0 + 2.0) \times 9.8$?}{No: only the cart presses on the table, $N = 49$ N.}
\WATCH{$T$ goes down with friction. Ask: does the string now pull harder or softer on the cart? (Harder: 16.8 N, it fights friction too.)} \tabularnewline[10pt]
\multicolumn{2}{@{}l}{\large\bfseries Frame 5: exit check and summary \hfill{\mono\small 38:00--50:00}}\tabularnewline*[2pt]
a. A 3.0 kg block on a frictionless table is tied over a pulley to a 1.5 kg hanging mass. Find the acceleration.\par
b. Find the tension in the string.\par
c. An Atwood machine holds 2.0 kg and 8.0 kg. Which side goes down, and how fast does it speed up?\par
\textbf{Summary.} Answer the essential question in two or three sentences. Use the words ``system'' and ``tension.'' &
\step{5.1}{black}{SAY}{\clock{38:00}Exit check on paper, alone, 7 minutes.}
\step{5.2}{black}{SAY}{\clock{45:00}Remove the covers: 3.3 m/s$^2$; 9.8 N; the 8.0 kg side, 5.9 m/s$^2$. 4 of 5 right on a and b: Day 5 (springs) as planned. Fewer: open Day 5 with one more table-and-hanging problem.}
\step{5.3}{black}{SAY}{\clock{45:00}Summary, 3 minutes. Sample: treat the tied objects as one system, where tension is internal and cancels, so $a$ = net external force $\div$ total mass; then look at one object alone to find the tension.}
\step{5.4}{black}{SAY}{\clock{48:00}Up next: ``Tomorrow: springs. A force you can see by how far something stretches.''}
\end{longtable}
\end{document}
